\bookchapter{الفصل 59: التعرف على iOS — بيئة Apple وبيئة التطوير}

\chapterintro{مبارك لك هذا القرار — أن تتحول من «مستخدم لـ iPhone» إلى «كاتب تطبيقات لـ iPhone». وقبل الشروع في العمل، لنفرش الخريطة: من أين جاء iOS؟ ماذا تعني «الحديقة المسوّرة» لـ Apple بالنسبة إلى المطورين؟ وماذا يجري بين أن تكتب الكود بأصابعك وأن يصل التطبيق إلى هاتف المستخدم؟}

\booksection{أولاً: نشأة iOS — هاتف واحد يغيّر صناعة كاملة}

في يناير 2007 أخرج جوبز أول iPhone في مؤتمر الإعلان. ولم يكن يدعم عندئذ حتى تطبيقات الأطراف الثالثة — حتى أطلقت Apple في 2008 متجر App Store وحزمة iPhone SDK وسمحت لمطوري العالم بكتابة البرامج لـ iPhone. في تلك اللحظة وُلد «اقتصاد التطبيقات»: حتى اليوم يضم App Store ملايين التطبيقات، وتجاوزت الحصص المدفوعة للمطورين تراكمياً مئات المليارات من الدولارات.

النظام نفسه يتطور أيضاً: سُمي في البداية iPhone OS ثم تغيّر اسمه إلى iOS في 2010؛ وتقدّمت الأرقام حتى iOS 17 (2023)، ثم غيّرت Apple في 2025 التسمية لتتبع السنة — أُطلق iOS 26 في 2025، والإصدار الحالي في أغسطس 2026 هو iOS 27. تتغير التسميات، لكن أمراً واحداً لم يتغير: عادة المستخدمين في الدفع وسمعة جودة البرمجيات في iOS هما الأفضل إجماعاً في بيئة الهواتف المتنقلة — التطبيق نفسه يدرّ على iOS دخلاً أعلى في الغالب بوضوح من نظيره على Android.

\begin{center}
\begin{tabularx}{\linewidth}{|X|X|X|}
\hline
\textbf{المحطة} & \textbf{السنة} & \textbf{ما يهم المطورين} \tabularnewline
\hline
إطلاق App Store & 2008 & انطلاق بيئة تطوير الأطراف الثالثة \tabularnewline
إطلاق Swift & 2014 & ظهور لغة تطوير iOS الحديثة \tabularnewline
إطلاق SwiftUI & 2019 & إطار واجهات تصريحي، وهو خط هذا الكتاب الرئيسي \tabularnewline
iOS 26 (العام الأول للتسمية الجديدة) & 2025 & التسمية تتبع السنة؛ ودمج شامل لـ Apple Intelligence \tabularnewline
iOS 27 & 2026 & أحدث إصدار حالي (مرجع هذا الكتاب) \tabularnewline
\hline
\end{tabularx}
\end{center}

\begin{center}\small\color{gray}محطات بارزة في تطوير iOS (حتى أغسطس 2026)\end{center}

\booksection{ثانياً: الحديقة المسوّرة — «قواعد» بيئة Apple و«أرباحها»}

توصف بيئة Apple كثيراً بأنها «حديقة مسوّرة» (Walled Garden) — قواعد كثيرة ورقابة صارمة: لا يُنشر التطبيق إلا بعد المراجعة، ولا يجوز التطوير إلا بـ Mac وXcode، والتوزيع يجري أساساً عبر App Store. يبدو ذلك خانقاً؟ لكن هذه القواعد جلبت للمطورين أرباحاً حقيقية ملموسة:

\begin{itemize}
  \item تفتت الأجهزة ضئيل جداً: طرازات iPhone معدودة، ومقاسات الشاشة وإصدارات النظام متجانسة للغاية — تكيّف مع 5 طرازات فتغطي 95\% من المستخدمين (بينما قد يتطلب Android التكيف مع عشرات الطرازات)؛
  \item نسبة تحديث إصدارات النظام عالية جداً: خلال أسابيع من إطلاق نظام جديد يرقّي إليه سبعون إلى ثمانون في المئة من المستخدمين، فيمكنك استخدام واجهات API الجديدة بجرأة دون مراعاة «نظام قديم عمره خمس سنوات»؛
  \item عادة الدفع جيدة: مستخدمو App Store أعدّ للإنفاق على التطبيقات والمشتريات الداخلية؛
  \item المراجعة صارمة لكن المتجر نظيف: القلّادات والبرامج الشريرة قليلة، فالمنتجات الجيدة أسهل في أن تُرى.
\end{itemize}

\begin{tcbwarn}{تنبيه: لنوضح التكلفة أولاً}
لتطوير iOS عقبة صلبة: تحتاج جهاز Mac (فـ Xcode حصري لـ macOS)، والنشر يتطلب حساب مطوّر بـ 99 دولاراً في السنة (حساب المطور الفردي المجاني يكفي للتصحيح على جهاز حقيقي محلياً لكنه لا يصلح للنشر). فكّر جيداً في هذا الاستثمار قبل الانطلاق — لكن الخبر السار أن الحساب المجاني يكفي تماماً في مرحلة التعلم، وسيفصّل ذلك الفصل 60.

\end{tcbwarn}

\booksection{ثالثاً: بأي لغة نكتب — Objective-C أم Swift؟}

في تاريخ تطوير iOS بطلا «رئيسيان» أيضاً. Objective-C هو الشيخ المؤسس — استُخدم من 2008 إلى 2014؛ وSwift هي اللغة الجديدة التي أطلقتها Apple في 2014، وتطورت بعدها حتى Swift 6 (أُطلقت في 2024 وأدخلت فحص تزامن صارماً)، وفي 2026 تستخدم Swift المشاريع الجديدة على وجه التقريب كلها.

كلمة صادقة لك وأنت في الصفر: تعلّم Swift مباشرة، ولا تنظر إلى Objective-C حتى — إلا إذا احتجت بعد التوظيف إلى صيانة كود تراثي من عشر سنوات. مزايا Swift: صياغة نقية (بلا أقواس منقوطة ولا أعباء ملفات الترويسة من عصر C)، وآلية القيم الاختيارية (Optional) التي أبادت «انهيار المؤشر الفارغ» من مستوى اللغة نفسها، وتكامل سلس مع SwiftUI، فضلاً عن Swift Playgrounds على iPad للتعلم واللعب معاً.

نمط Objective-C (أطلال قديمة)

\begin{lstlisting}
NSString *name = @"أحمد"; [name stringByAppendingString:@"مرحبا"]; // الأقواس المعقوفة والمؤشرات والفواصل المنقوطة — ثلاثية كاملة
\end{lstlisting}

نمط Swift (عصري ونقي)

\begin{lstlisting}
var name = "أحمد" name += "مرحبا" // واضح ومباشر كاللغة الطبيعية
\end{lstlisting}

\booksection{رابعاً: كيف يولد التطبيق؟}

ككل المنصات، لولادة تطبيق iOS خط إنتاج أيضاً. لكن نسخة Apple فيها «فحصان» إضافيان:

فكرة

التصميم (الواجهة + الوظائف)

كتابة الكود (Swift + SwiftUI)

البناء والتوقيع

الاختبار (محاكي + جهاز حقيقي)

مراجعة Apple

النشر في App Store

\begin{center}\small\color{gray}الشكل 1: المسار الكامل لتطبيق iOS من الفكرة إلى النشر (لاحظ مرحلة «المراجعة» الإضافية)\end{center}

مفاهيم أساسية في المسار:

\begin{itemize}
  \item التوقيع (Signing): كل تطبيق iOS يجب أن يُوقَّع ليُثبَّت — إثباتاً بأن «هذه الحزمة منك ولم يُعبث بها». في مرحلة التعلم يوقّع Xcode تلقائياً بمعرّف Apple ID الخاص بك، وعند النشر تُستخدم شهادة المطور. التفصيل في الفصل 67.
  \item المحاكي (Simulator): جهاز iPhone افتراضي على Mac وساحة العمل الرئيسية للتطوير والتصحيح. يعمل بمعالج Mac نفسه، لذا فهو سلس للغاية وتجربته أفضل بكثير من محاكي Android.
  \item TestFlight: أداة Apple الرسمية «لتوزيع الاختبار الداخلي» — تُرسل قبل النشر نسخة تجريبية إلى ما يصل إلى 10000 مستخدم اختبار يثبّتونها ويجربونها عبر تطبيق TestFlight. وهذه أناقة يفرد بها iOS وحده.
\end{itemize}

\booksection{خامساً: منظومة تقنيات iOS الكاملة — «قائمة أسلحتك»}

\begin{center}
\begin{tabularx}{\linewidth}{|X|X|X|}
\hline
\textbf{الطبقة} & \textbf{الخيار السائد (2026)} & \textbf{ما استخدامه} \tabularnewline
\hline
اللغة & Swift 6 & لغة كتابة المنطق \tabularnewline
الواجهة & SwiftUI + دليل تصميم Apple (HIG) & رسم الواجهة بالكود (خط هذا الكتاب الرئيسي) \tabularnewline
البنية & MVVM + @Observable & فصل الواجهة عن المنطق \tabularnewline
تخزين البيانات & SwiftData (قاعدة البيانات) + UserDefaults (الإعدادات) & حفظ البيانات في الهاتف \tabularnewline
الشبكة & URLSession + async/await + Codable & التواصل مع الخادم \tabularnewline
أداة البناء & Xcode Build System / xcodebuild & تحويل الكود إلى حزمة تثبيت \tabularnewline
بيئة التطوير & Xcode (حصري لـ macOS) & بيئة التطوير المتكاملة الرسمية من Apple \tabularnewline
\hline
\end{tabularx}
\end{center}

\begin{center}\small\color{gray}نظرة سريعة على منظومة تقنيات iOS السائدة في 2026\end{center}

هل لاحظت؟ هذه القائمة أقصر من قائمة Android — لأن Apple «تكفل» بأمور كثيرة. ثمن البيئة المغلقة هو حرية أقل، ومكافأتها راحة البال: لا تتحير في اختيار مكتبة الشبكة أو قاعدة البيانات، فما تقدمه الجهة الرسمية هو غالباً الحل الأمثل.

\begin{tcbpractice}{لنطبّق: ألقِ نظرة على «هوية Apple» في هاتفك}
افتح في iPhone [الإعدادات] [عام] [حول]، وسجّل ثلاثة أمور: إصدار iOS واسم الطراز والمساحة المتاحة. ثم تجول على صفحة «اليوم» في App Store وراقب لقطات شاشة التطبيقات المميزة ونصوصها — فمن اليوم يتغير منظرك في المتجر من «مستخدم» إلى «مطور».

\end{tcbpractice}

\begin{tcbsummary}{خلاصة الفصل}
\begin{itemize}
  \item في 2008 أُطلق App Store ووُلد «اقتصاد التطبيقات»؛ والإصدار الحالي في 2026 هو iOS 27
  \item حساب الحديقة المسوّرة: قواعد كثيرة، لكن أجهزة قليلة وإصدارات جديدة ومستخدمون يبذلون المال ومتجر نظيف
  \item العقبة الصلبة: Mac ضروري؛ والنشر يتطلب حساب مطوّر بـ 99 دولاراً سنوياً؛ والحساب المجاني يكفي في مرحلة التعلم
  \item تعلّم Swift مباشرة؛ فـ Objective-C لا يخص إلا الكود التراثي
  \item خط الإنتاج أطول بخطوة من Android: مراجعة Apple؛ وTestFlight أداة توزيع الاختبار الداخلي المذهلة
  \item منظومة 2026 القياسية: Swift 6 + SwiftUI + MVVM + SwiftData + URLSession + Xcode
\end{itemize}

\end{tcbsummary}
